% TEMPLATE-MILC-v3.tex - plantilla maestra UNICA de las guias MILC v3.
%
% La usan las tres familias de guias, cada una con su driver:
%   · Grados 6-11  -> scripts/build-guias-g11.py
%   · Bebras       -> scripts/build-guias-bebras.py
%   · Semillero    -> scripts/build-guias-semillero.py
%
% Antes habia tres copias casi identicas (~400 lineas iguales, 6 distintas):
% cada mejora habia que aplicarla tres veces, y en la practica no se hacia
% (el motor de recursos vivia solo en dos de ellas). Lo que varia entre
% familias esta parametrizado en 6 placeholders que cada driver llena
% (se nombran aqui SIN su sintaxis real: el builder reemplaza por texto plano,
% tambien dentro de los comentarios, y un valor multilinea rompe el preambulo):
%
%   HEADER_IZQ        encabezado izquierdo de pagina
%   HEADER_DER        encabezado derecho de pagina
%   PORTADA_ETIQUETA  rotulo sobre el numero grande (GRADO/MOMENTO/MODULO)
%   PORTADA_SERIE     linea de serie de la portada
%   PORTADA_ID        identificador de la guia en la portada
%   PORTADA_CIRCULO   el circulito de grado (°), o vacio si no aplica
%
% El resto de claves las documenta content/guias/_SCHEMA.md. El script valida
% que no queden placeholders sin reemplazar antes de escribir el archivo final.

\documentclass[11pt,letterpaper]{article}
\usepackage[margin=2.0cm]{geometry}
\usepackage[table]{xcolor}
\usepackage{fontspec}
\usepackage[spanish]{babel}
\usepackage{tikz}
% Librerías TikZ para los diagramas de las guías (bloque `recursos:`):
%  · babel  → babel en español vuelve activos caracteres como «>», y TikZ los usa
%             en sus opciones (>=stealth, ->). Sin esta librería, cualquier
%             diagrama con flechas revienta con «\language@active@arg> has an
%             extra }». Debe ir ANTES de que se dibuje nada.
%  · positioning        → sintaxis «below left=2pt and 3pt».
%  · decorations.pathreplacing → llaves ({brace}) para señalar rangos.
\usetikzlibrary{babel,positioning,decorations.pathreplacing}
\usepackage{graphicx}
% Circuitos: el trabajo de electrónica se hace hoy con micro:bit y sus sensores
% integrados (MakeCode, sin cableado), así que no se necesitan esquemáticos.
% Si algún día entran montajes con protoboard, basta descomentar esta línea:
% los diagramas .tex/.tikz declarados en `recursos:` ya se incrustan solos.
% \usepackage{circuitikz}
\usepackage[most]{tcolorbox}
\usepackage{tabularx}
\usepackage{array}
\usepackage{fancyhdr}
\usepackage{titlesec}
\usepackage[document]{ragged2e}  % bandera a la derecha en todo el cuerpo
\usepackage{enumitem}
\usepackage{microtype}

% Helvetica Neue no trae la flecha U+2192, asi que cada «→» escrito literalmente
% en el YAML se compilaba a NADA, en silencio (462 flechas en 61 guias: rutas del
% tipo «paleta → tipografia → lockup» salian rotas). Mapeamos el caracter a su
% version matematica, que si tiene glifo. Asi el autor puede seguir escribiendo
% «→» comodamente en el YAML.
\usepackage{newunicodechar}
\newunicodechar{→}{\ensuremath{\rightarrow}}
% Mismo problema con los simbolos de marca y vinetas: el «✓» de «una palomita ✓»
% salia como cuadrito vacio en el PDF de todas las guias que lo usaban.
\usepackage{pifont}
\newunicodechar{✓}{\ding{51}}
\newunicodechar{✗}{\ding{55}}
\newunicodechar{✔}{\ding{52}}
\newunicodechar{•}{\textbullet}
\newunicodechar{≥}{\ensuremath{\geq}}
\newunicodechar{≤}{\ensuremath{\leq}}

\setmainfont{Helvetica Neue}
\setsansfont{Helvetica Neue}
\setlength{\parindent}{0pt}
\setlength{\parskip}{6pt}
% Sin particion de palabras: en 6.º-7.º una palabra cortada por guion al final
% de linea («Comuni-cacion») frena la lectura mucho mas que una linea corta.
\hyphenpenalty=10000
\exhyphenpenalty=10000
\pretolerance=2000
\tolerance=3000
\setlength{\emergencystretch}{3em}
\linespread{1.10}
\renewcommand{\arraystretch}{1.25}
\setlist{leftmargin=5mm,itemsep=1mm,topsep=1mm}
\newcolumntype{Y}{>{\RaggedRight\arraybackslash}X}

\definecolor{milcMagenta}{HTML}{E6007E}
\definecolor{milcVino}{HTML}{5A0038}
\definecolor{milcTurquesa}{HTML}{0093A5}
\definecolor{milcVerde}{HTML}{58A923}
\definecolor{milcMostaza}{HTML}{E5B400}
\definecolor{milcOcre}{HTML}{B86600}
\definecolor{milcNegro}{HTML}{241816}
\definecolor{milcGris}{HTML}{F7F7F4}
\definecolor{milcLinea}{HTML}{E8E2DD}
\definecolor{milcGradePrimary}{HTML}{84CC16}
\definecolor{milcGradeDark}{HTML}{4D7C0F}
\definecolor{milcGradeSoft}{HTML}{ECFCCB}
\definecolor{milcGradeAccent}{HTML}{CFE7FF}
\definecolor{milcGradeText}{HTML}{FFFFFF}
\color{milcNegro}

\pagestyle{fancy}
\fancyhf{}
\lhead{\small Tecnología e Informática · Grado 7}
\rhead{\small Guía 8 · MILC v3}
\cfoot{\small\thepage}
\renewcommand{\headrulewidth}{0pt}

\newtcolorbox{titlebox}[1]{
  enhanced,colback=#1,colframe=#1,arc=7mm,boxrule=0pt,
  left=7mm,right=7mm,top=3mm,bottom=3mm,
  before skip=8pt,after skip=8pt,
  fontupper=\sffamily\bfseries\Large\color{white}
}

\newtcolorbox{softbox}[3]{
  enhanced,colback=#2,colframe=#1,borderline west={4pt}{0pt}{#1},
  arc=2mm,boxrule=.4pt,left=4mm,right=4mm,top=3mm,bottom=3mm,
  before skip=6pt,after skip=8pt,title={#3},coltitle=white,
  colbacktitle=#1,fonttitle=\sffamily\bfseries,fontupper=\RaggedRight,
  attach boxed title to top left={xshift=3mm,yshift=-2mm},
  boxed title style={arc=2mm, boxrule=0pt}
}

\newtcolorbox{drawbox}[2]{
  enhanced,colback=white,colframe=#1,arc=4mm,boxrule=1pt,
  left=5mm,right=5mm,top=4mm,bottom=4mm,before skip=8pt,after skip=8pt,
  title={#2},coltitle=white,colbacktitle=#1,fonttitle=\sffamily\bfseries,
  fontupper=\RaggedRight,
  attach boxed title to top left={xshift=4mm,yshift=-2mm},
  boxed title style={arc=3mm, boxrule=0pt}
}

\newtcolorbox{infoband}[1]{
  enhanced,colback=#1!8,colframe=#1!50,arc=2mm,boxrule=.5pt,
  left=4mm,right=4mm,top=2mm,bottom=2mm,before skip=4pt,after skip=4pt,
  fontupper=\small\RaggedRight
}

% Puente narrativo entre secciones (contrato editorial Conéctate).
% Da continuidad: "Acabas de X. Ahora vas a Y porque Z."
% Visualmente: banda izquierda en color del grado, texto en itálica pequeña.
\newtcolorbox{puentebox}{
  enhanced,colback=milcGradeSoft,colframe=milcGradeDark,
  borderline west={3pt}{0pt}{milcGradeDark},
  arc=0mm,boxrule=0pt,left=5mm,right=4mm,top=2.5mm,bottom=2.5mm,
  before skip=4pt,after skip=8pt,
  fontupper=\itshape\small\color{milcNegro}\RaggedRight
}
% La flecha va en modo matematico a proposito: Helvetica Neue NO tiene el glifo
% U+2192, asi que \textrightarrow se compilaba a NADA («Missing character: There
% is no → in font Helvetica Neue») y el puente salia sin flecha. En modo
% matematico la flecha viene de la fuente matematica y si se dibuja.
\newcommand{\puente}[1]{%
  \begin{puentebox}%
    \textcolor{milcGradeDark}{$\rightarrow$}\hspace{2mm}#1%
  \end{puentebox}%
}

\newcommand{\linead}[1]{\noindent\color{milcLinea}\rule{#1}{0.5pt}\color{milcNegro}}
\newcommand{\respuesta}[1]{\par\vspace{2mm}\foreach \n in {1,...,#1}{\noindent\color{milcLinea}\rule{\linewidth}{0.5pt}\par\vspace{5mm}}\color{milcNegro}}
\newcommand{\checkbox}{\textcolor{milcGradeDark}{\fbox{\phantom{x}}}\hspace{5pt}}

\newcommand{\phasecard}[4]{
\begin{tcolorbox}[enhanced,colback=#3,colframe=#2,arc=3mm,boxrule=.5pt,height=3.05cm,valign=center,halign=center,left=1.5mm,right=1.5mm,top=2mm,bottom=2mm]
{\sffamily\bfseries\fontsize{22}{24}\selectfont\textcolor{#2}{#1}}\par
{\sffamily\bfseries\small #4}
\end{tcolorbox}
}

% ── Piezas del rediseno 2026 (legibilidad 6.º-11.º) ───────────────────
% Banda de estacion: reemplaza el titlebox macizo. Titulo a la izquierda,
% dato practico (tiempo/modalidad) a la derecha, en una sola linea.
\newcommand{\milcestacion}[2]{%
\begin{tcolorbox}[enhanced,colback=#1,colframe=#1,arc=2mm,boxrule=0pt,
  left=5mm,right=5mm,top=2.6mm,bottom=2.6mm,before skip=11pt,after skip=7pt]
{\sffamily\bfseries\fontsize{15}{18}\selectfont\color{white}#2}
\end{tcolorbox}}

% Tarjeta de paso numerada: sustituye las tablas «Pilar / Aplicacion», que
% eran formato de consulta docente, no de estudiante. El numero va en un
% circulo sobre el borde para que la secuencia se lea de un vistazo.
\newcommand{\milcpaso}[3]{%
\begin{tcolorbox}[enhanced,colback=white,colframe=milcLinea,arc=1.5mm,boxrule=.9pt,
  left=14mm,right=4mm,top=3mm,bottom=3mm,before skip=4pt,after skip=4pt,
  fontupper=\RaggedRight,
  overlay={\node[anchor=north west,circle,fill=#2,text=white,inner sep=0pt,
    minimum size=7.4mm,font=\sffamily\bfseries\small]
    at ([xshift=3.6mm,yshift=-3.2mm]frame.north west) {#1};}]
#3
\end{tcolorbox}}

% Casilla marcable de tamano util con lapiz (la anterior era \fbox{\phantom{x}},
% de alto variable segun el texto que la seguia).
\newcommand{\milcmarca}{\framebox[3.8mm]{\rule{0pt}{3.4mm}}}

% Fila marcable de lista suelta (checks de la estacion 1).
\newcommand{\milccheckrow}[1]{%
\par\vspace{1.8mm}\noindent
\makebox[6.4mm][l]{\raisebox{-.55mm}{\textcolor{milcNegro}{\milcmarca}}}%
\parbox[t]{\dimexpr\linewidth-6.4mm\relax}{\RaggedRight #1}\par}

% Celda marcable dentro de la rubrica (tabla de autoevaluacion). Misma
% sangria francesa, pero pensada para vivir dentro de una columna X.
\newcommand{\milcopcion}[1]{%
\makebox[5.2mm][l]{\raisebox{-.55mm}{\textcolor{milcNegro}{\milcmarca}}}%
\parbox[t]{\dimexpr\linewidth-5.2mm\relax}{\RaggedRight #1}}

% ── Recursos visuales de la guía ──────────────────────────────────────
% Declarados en el bloque `recursos:` del YAML y emitidos por el builder.
% \guiaFigura   → imagen rasterizada o PDF (foto, carta celeste, montaje).
% \guiaDiagrama → fuente TikZ/circuitikz (.tex) incrustada como vector:
%                 es la vía para circuitos y esquemáticos editables.
% #1 = ruta relativa al .tex · #2 = pie (puede ir vacío)
\newcommand{\guiaFigura}[2]{%
\begin{center}
\includegraphics[width=0.88\linewidth,height=0.38\textheight,keepaspectratio]{#1}\\[1.5mm]
{\footnotesize\itshape\textcolor{milcNegro!75}{#2}}
\end{center}
}

\newcommand{\guiaDiagrama}[2]{%
\begin{center}
\input{#1}\\[1.5mm]
{\footnotesize\itshape\textcolor{milcNegro!75}{#2}}
\end{center}
}

\begin{document}
\thispagestyle{empty}

% ── Portada ───────────────────────────────────────────────────────────
% Antes: un rectangulo de color a pagina completa. Se veia bien en pantalla,
% pero en la fotocopiadora del colegio se comia el toner y salia gris sucio.
% Ahora la banda ocupa el tercio superior y el resto va en blanco.
\begin{tikzpicture}[remember picture,overlay]
  \path[fill=milcGradePrimary]
    (current page.north west) rectangle ([yshift=-10.6cm]current page.north east);

  \node[anchor=north west,text=milcGradeText] at ([xshift=2.0cm,yshift=-1.55cm]current page.north west)
    {\sffamily\bfseries\fontsize{12}{14}\selectfont GRADO};
  \node[anchor=north west,text=milcGradeText,opacity=.85] at ([xshift=2.0cm,yshift=-2.35cm]current page.north west)
    {\sffamily\bfseries\fontsize{8.5}{10}\selectfont SERIE GUÍAS · TECNOLOGÍA E INFORMÁTICA};
  \node[anchor=north east,text=milcGradeText,opacity=.85,align=right] at ([xshift=-2.0cm,yshift=-1.55cm]current page.north east)
    {\sffamily\bfseries\fontsize{9}{12}\selectfont I.E. Sor María Juliana\\Cartago, Valle del Cauca};

  \node[anchor=west,text=milcGradeText] (gradonum) at ([xshift=1.85cm,yshift=-7.1cm]current page.north west)
    {\sffamily\bfseries\fontsize{112}{112}\selectfont 7};
  % El circulito de grado (°) solo aplica a las guias de grado («11°»); en
  % Bebras y Semillero el numero grande es un momento/modulo, no un grado.
  % Se posiciona relativo al nodo `gradonum`, no en coordenadas absolutas.
  \draw[milcGradeText,line width=1.6pt]
    ([xshift=2.0mm,yshift=-4.5mm]gradonum.north east) circle (.15cm);

  \node[anchor=west,text=milcGradeText,text width=11.4cm,align=left] at ([xshift=1.5cm,yshift=0.5cm]gradonum.east)
    {
     {\sffamily\bfseries\fontsize{9}{11}\selectfont\MakeUppercase{Período 1 · Trabajo colaborativo en la nube}}\\[3mm]
     {\sffamily\bfseries\fontsize{21}{25}\selectfont\setlength{\baselineskip}{27pt}Versiones e\\[4pt]historial ---\\[4pt]el sistema\\[4pt]recuerda todo\\[4pt]lo que hice}\\[3.5mm]
     {\sffamily\bfseries\fontsize{15}{18}\selectfont Guía 8-7-TIC}
    };
\end{tikzpicture}

\vspace*{9.4cm}
\vfill

{\sffamily\bfseries\fontsize{8}{10}\selectfont\textcolor{milcGradeDark}{LO QUE TE LLEVAS HOY}}

\begin{tcolorbox}[enhanced,colback=white,colframe=milcNegro,arc=2mm,boxrule=1.6pt,
  left=5mm,right=5mm,top=4mm,bottom=4mm,before skip=3pt,after skip=12pt,
  fontupper=\RaggedRight\fontsize{11}{15.5}\selectfont]
Una \textbf{auditoría del historial} de uno de tus archivos en OneDrive: \textbf{(a) abres el historial de versiones} de un documento; \textbf{(b) identificas} mínimo 3 versiones distintas con sus fechas; \textbf{(c) practicas restaurar} una versión anterior y volver a la actual; \textbf{(d) en el cuaderno}: dibujo de la pantalla \emph{Historial de versiones} con las versiones identificadas + las 4 reglas de uso responsable del historial.
\end{tcolorbox}

\begin{tcolorbox}[enhanced,colback=milcGris,colframe=milcGris,arc=2mm,boxrule=0pt,
  left=5mm,right=5mm,top=3.5mm,bottom=3.5mm,after skip=12pt]
{\sffamily\bfseries\fontsize{8}{10}\selectfont\textcolor{milcNegro!55}{ESTA GUÍA ES DE}}\\[3mm]
\begin{tabularx}{\linewidth}{X p{3.2cm} p{3.2cm}}
\linead{\linewidth} & \linead{3.0cm} & \linead{3.0cm}\\[-1mm]
{\fontsize{8.5}{10}\selectfont\textcolor{milcNegro!55}{Nombre completo}} &
{\fontsize{8.5}{10}\selectfont\textcolor{milcNegro!55}{Grupo}} &
{\fontsize{8.5}{10}\selectfont\textcolor{milcNegro!55}{Fecha}}\\
\end{tabularx}
\end{tcolorbox}

\vfill

\noindent\textcolor{milcLinea}{\rule{\linewidth}{1pt}}\\[2mm]
{\sffamily\bfseries\fontsize{9.5}{12}\selectfont Autor: Dr. Álvaro Cárdenas Orozco}\\[1mm]
{\sffamily\fontsize{8.5}{11}\selectfont\textcolor{milcNegro!55}{Metodología MILC · apertura ancestral · pensamiento computacional · triángulo Dussel–estoicismo–Floridi}}

\newpage

\section*{Apertura: Versiones e historial --- el sistema recuerda todo lo que escribiste}

\begin{softbox}{milcGradeDark}{milcGradeSoft}{Saber ancestral}
\textbf{En el taller de doña Carmen la costurera del Valle del Cauca, había una costumbre estricta: nunca botar las piezas viejas.} Cuando un cliente le pedía un vestido a doña Carmen y después del primer corte cambiaba de idea (\emph{``mejor más largo''}, \emph{``mejor mangas cortas''}), \textbf{ella guardaba TODOS los pedazos cortados}. Si después al cliente se le ocurría \emph{``no, mejor sí lo quería corto''}, doña Carmen no tenía que empezar de cero: \emph{``Aquí está el pedazo, lo volvemos a meter''}. \textbf{Cada decisión tenía marcha atrás}. Esa virtud --- guardar lo que se ha hecho para poder regresar a versiones anteriores --- es exactamente lo que hace OneDrive con tus archivos. Cada vez que guardas un cambio, \textbf{no se borra el anterior}: queda \emph{una versión nueva}, y la anterior queda guardada por si te arrepientes. \emph{En el campo se llamaba ``no botar las piezas''. En digital se llama ``historial de versiones''.}
\end{softbox}

\begin{softbox}{milcTurquesa}{milcGris}{Saber contemporáneo}
El \textbf{historial de versiones} es una función automática de OneDrive (y de Google Drive, y de la mayoría de servicios en la nube). Cada vez que el archivo se modifica de manera significativa --- o cada cierto tiempo --- \textbf{el sistema guarda una versión completa} del archivo. Tú puedes ver el historial de versiones, leer cualquier versión anterior, e incluso \textbf{restaurar} una versión vieja como la actual. \textbf{¿Para qué sirve?} (1) Si cambias un archivo y te arrepientes, vuelves a la versión anterior. (2) Si un compañero (con permiso editor) daña algo, restauras tu versión. (3) Si quieres saber qué decía el documento ayer, lo abres en versión anterior sin cambiar el actual. (4) Si pierdes algo por error, no se pierde para siempre. \textbf{En Microsoft 365, OneDrive guarda hasta 25 versiones de cada archivo durante 30 días}. Algunas empresas configuran más.
\end{softbox}

\begin{softbox}{milcMostaza}{milcGris}{Pregunta-puente}
\textit{Estuviste 2 horas escribiendo un trabajo en línea. De repente borras por error un párrafo importante y guardas sin querer. ¿\textbf{Se perdió para siempre} o hay forma de recuperarlo sin haber hecho copia manual?}
\end{softbox}

\begin{softbox}{milcVerde}{milcGris}{Saber y saber hacer}
Al terminar la clase: (1) podrás \textbf{identificar} qué es el historial de versiones y cómo se accede; (2) sabrás \textbf{aplicar} la restauración de una versión anterior; (3) podrás \textbf{evaluar} cuándo usar el historial vs hacer copia manual; (4) habrás \textbf{auditado} el historial de uno de tus archivos.
\end{softbox}



\small
\begin{tabularx}{\textwidth}{>{\bfseries}p{3.0cm}Y}
\rowcolor{milcGradePrimary!12}Autor & Dr. Álvaro Cárdenas Orozco\\
DBA & Utiliza herramientas digitales para trabajar de manera colaborativa con otros (MEN, T\&I 7°)\\
Referentes & Saber ancestral de la minga · Microsoft 365 y OneDrive · Coautoría asincrónica\\
Producto final & Una \textbf{auditoría del historial} de uno de tus archivos en OneDrive: \textbf{(a) abres el historial de versiones} de un documento; \textbf{(b) identificas} mínimo 3 versiones distintas con sus fechas; \textbf{(c) practicas restaurar} una versión anterior y volver a la actual; \textbf{(d) en el cuaderno}: dibujo de la pantalla \emph{Historial de versiones} con las versiones identificadas + las 4 reglas de uso responsable del historial.\\
\end{tabularx}
\normalsize

\puente{Hoy haces 4 cosas: imaginas el caso de pérdida, aprendes a abrir el historial, restauras una versión, auditas tu archivo.}

\milcestacion{milcGradeDark}{Ruta MILC: así vamos a aprender}
\begin{tabularx}{\textwidth}{>{\centering\arraybackslash}X>{\centering\arraybackslash}X>{\centering\arraybackslash}X>{\centering\arraybackslash}X}
\phasecard{1}{milcVerde}{milcGris}{Escucha\\[-1mm]\small Observo, escucho y pregunto.} &
\phasecard{2}{milcTurquesa}{milcGris}{Entender\\[-1mm]\small Sistematizo y ordeno.} &
\phasecard{3}{milcMagenta}{milcGris}{Hacer\\[-1mm]\small Construyo y aplico.} &
\phasecard{4}{milcVino}{milcGris}{Cierre\\[-1mm]\small Evalúo y reflexiono.}
\end{tabularx}


\puente{Empezamos con 4 casos comunes de pérdida y cómo se solucionan.}

\milcestacion{milcVerde}{Estación 1 de 3 · Escucha}

\begin{softbox}{milcVerde}{milcGris}{Escena para escuchar}
\textbf{Actividad 1 · ANALIZA --- 4 casos comunes de pérdida} (10 min · individual). Lee estos 4 casos. Para cada uno, escribe en el cuaderno \textbf{si el historial de versiones podría salvar a la persona y por qué}. \emph{(1) Lucía escribió 3 párrafos en su Word, decidió empezar de cero, los borró y guardó. Una hora después se arrepintió.} \emph{(2) Pedro compartió su archivo con un compañero como editor. El compañero, sin querer, borró toda la sección de Pedro. (3) María entregó su trabajo al profe. El profe lo modificó por error y le aceptó cambios. María quiere ver el documento como estaba antes de la modificación. (4) David trabajó en un Excel durante 1 hora. La sesión expira por inactividad y al volver ve cosas raras. Quiere saber qué cambió.}
\end{softbox}

{\sffamily\bfseries\fontsize{8}{10}\selectfont\textcolor{milcNegro!55}{MARCA CUANDO LO HAYAS HECHO}}
\milccheckrow{Leí los 4 casos con calma.}
\milccheckrow{Decidí si el historial ayudaría en cada uno.}
\milccheckrow{Pensé qué se haría sin historial.}

\begin{infoband}{milcVerde}


\textbf{Lo que va al cuaderno (Actividad 1 · ANALIZA).} Título: «Actividad 1 --- 4 casos de pérdida». \textbf{Formato:} lista numerada con respuesta corta. \textbf{Extensión:} 4 líneas. \textbf{Sabes que terminaste cuando} reconoces que el historial salva en los 4 casos.
\end{infoband}

\puente{Después aprendes cómo abrir y leer el historial, y las 4 reglas de uso responsable.}

\milcestacion{milcTurquesa}{Estación 2 de 3 · Entender}

\begin{softbox}{milcTurquesa}{milcGris}{Lo que vas a entender}
Los 4 casos se solucionan con \textbf{historial de versiones}. Ahora aprende cómo acceder a él y las 4 reglas para usarlo bien.
\end{softbox}

\milcpaso{1}{milcTurquesa}{\textbf{Cómo abrir el historial de versiones en OneDrive (Word, Excel, PowerPoint).} \textbf{Método 1 desde el archivo abierto:} con el documento abierto, ve a la pestaña \emph{Archivo} (arriba a la izquierda) → \emph{Información} → \emph{Historial de versiones} (o \emph{Versiones anteriores}). Aparece un panel lateral con la lista de versiones (fecha y autor de cada cambio). \textbf{Método 2 desde OneDrive:} en OneDrive, encuentra el archivo, haz clic derecho → \emph{Historial de versiones}. Mismo panel. \textbf{Lo que ves:} lista de versiones, cada una con \emph{fecha y hora exactas} + \emph{nombre del autor que hizo el cambio}.}
\milcpaso{2}{milcTurquesa}{\textbf{Cómo navegar y restaurar versiones.} \textbf{(a) Ver una versión anterior:} haz clic en la fecha de la versión que te interesa. Se abre la versión en otra ventana o pestaña, solo lectura. La actual sigue intacta. \textbf{(b) Restaurar versión anterior:} si decides que esa versión es la que quieres mantener, hay un botón \emph{``Restaurar''}. Eso convierte esa versión vieja en la actual (la actual pasa a versión anterior). \textbf{(c) Copiar parte de versión anterior:} si solo quieres recuperar un párrafo específico, copia el texto de la versión anterior y pégalo en la actual.}
\milcpaso{3}{milcTurquesa}{\textbf{Las 4 reglas del uso responsable del historial.} (1) \textbf{Antes de restaurar, lee con calma.} A veces el problema no era la versión, sino tu memoria. Lee la versión vieja en otra pestaña antes de decidir restaurarla. (2) \textbf{Guarda nota de qué restauraste y por qué.} En el cuaderno o en un comentario del archivo. Si después aparece otro problema, sabes qué hiciste. (3) \textbf{No restaures versiones de hace mucho tiempo sin pensar.} Si el archivo lleva 1 mes evolucionando, restaurar la versión de hace 3 semanas pierde el trabajo de las últimas semanas. (4) \textbf{El historial no reemplaza el backup manual de cosas críticas.} Para trabajos importantes (proyecto final, tesis), también copia a USB o a otro lugar. Es como tener doble seguro.}
\milcpaso{4}{milcTurquesa}{\textbf{Diferencia entre historial vs papelera.} \textbf{Historial} = versiones de un archivo que aún existe (cambios dentro del archivo). \textbf{Papelera} = archivos completos que fueron borrados pero pueden recuperarse durante 30 días. Si borraste un archivo completo, usa la papelera. Si dañaste un archivo (pero existe), usa el historial. Las dos son herramientas distintas que se complementan.}

\begin{drawbox}{milcTurquesa}{Acceso + restauración + 4 reglas}
\textbf{Acceso:} Archivo → Información → Historial de versiones (o clic derecho desde OneDrive). \\[1mm]
\textbf{Restauración:} clic en fecha → leer versión → botón Restaurar. \\[1mm]
\textbf{4 reglas:} lee antes · guarda nota · no restaures muy atrás sin pensar · backup manual para lo crítico. \\[1mm]
\textbf{Historial vs Papelera:} historial cambia dentro del archivo · papelera recupera archivos borrados.
\end{drawbox}

\begin{softbox}{milcMostaza}{milcGris}{Errores comunes y su corrección}
\textbf{Mitos sobre el historial de versiones.} (1) \emph{``El historial guarda para siempre''} --- Falso. OneDrive guarda 25 versiones durante 30 días. Después se borran las viejas. (2) \emph{``Restaurar una versión borra el resto del historial''} --- Falso. La versión actual se vuelve una versión más; las viejas siguen. (3) \emph{``Solo el dueño ve el historial''} --- Falso. Cualquier persona con permiso editor puede verlo. (4) \emph{``El historial no funciona si trabajaste offline''} --- Cierto en parte: si trabajaste offline, los cambios se suben cuando vuelves a tener internet, pero las versiones se cuentan distinto.
\end{softbox}

\begin{infoband}{milcTurquesa}


\textbf{Lo que va al cuaderno (Actividad 2 · EXPLICA).} Título: «Actividad 2 --- Cómo usar el historial». \textbf{Formato:} bloque A con el camino paso a paso para abrir el historial. Bloque B con las 4 reglas. \textbf{Extensión:} 8 líneas. \textbf{Sabes que terminaste cuando} podrías guiar a un compañero a abrir el historial en su computador.
\end{infoband}

\puente{Con todo claro, abres uno de tus archivos y exploras su historial en OneDrive.}

\milcestacion{milcMagenta}{Estación 3 de 3 · Hacer}

\begin{softbox}{milcMagenta}{milcGris}{Cómo vas a construir tu producto}
\textbf{Actividad 3 · EVALÚA --- Auditoría del historial de uno de tus archivos} (28 min · individual). Vas a abrir uno de tus archivos en OneDrive y explorar su historial real. La idea es \emph{hacerte amigo} del historial antes de necesitarlo en una emergencia.
\end{softbox}

\milcpaso{1}{milcMagenta}{\textbf{Paso 1 --- Escoge el archivo.} Elige el archivo más activo que tengas: el de la S4 (mi-vision-minga-digital), el de la S5 (compartido), o el colaborativo de S6. Mientras más cambios haya tenido, más versiones tendrá su historial.}
\milcpaso{2}{milcMagenta}{\textbf{Paso 2 --- Abre el archivo y entra al historial.} En OneDrive: clic derecho sobre el archivo → \emph{Historial de versiones}. O dentro del archivo abierto: \emph{Archivo → Información → Historial de versiones}. Aparece un panel lateral con la lista de versiones.}
\milcpaso{3}{milcMagenta}{\textbf{Paso 3 --- Identifica 3 versiones.} En tu cuaderno escribe la lista de versiones que ves, con fecha exacta y autor: \emph{(1) Versión más reciente: [fecha] [tu nombre]. (2) Versión intermedia: [fecha] [autor]. (3) Versión más antigua que aparece: [fecha] [autor]}. Si hay menos de 3 versiones (el archivo es muy nuevo), edita el archivo, guarda, y vuelve al historial. Una vez tengas mínimo 3, sigue.}
\milcpaso{4}{milcMagenta}{\textbf{Paso 4 --- Abre una versión anterior y compárala con la actual.} Haz clic en la fecha de la versión intermedia (no la más antigua, una del medio). Se abre en otra pestaña. Compara: ¿qué tenía esa versión que no tiene la actual? ¿Qué se agregó? ¿Qué se quitó? Anota 2-3 diferencias en el cuaderno.}
\milcpaso{5}{milcMagenta}{\textbf{Paso 5 --- Practica restaurar (con cuidado).} Aquí lo importante: vas a restaurar la versión intermedia y después volver a la actual. \emph{(a)} En la versión intermedia abierta, busca el botón \textbf{``Restaurar''}. Haz clic. Acepta el mensaje. Ahora esa versión es la actual. \emph{(b) Inmediatamente} vuelve al historial: ahora la que era la actual antes está como una versión más arriba. \emph{(c) Haz clic en ella y restáurala}. Vuelves al estado original. Esto te entrena en que \textbf{restaurar no destruye, solo reordena}.}
\milcpaso{6}{milcMagenta}{\textbf{Paso 6 --- Dibujo del historial + 4 reglas firmadas + reflexión.} En el cuaderno haz un \textbf{boceto del panel de historial} (un panel lateral con lista de fechas y autores). Etiqueta: \emph{``Aquí veo todas las versiones de mi archivo''}. Después escribe las 4 reglas del uso responsable y palomita ✓ las que te comprometes a cumplir. Reflexión: \emph{``Hoy descubrí que en OneDrive nada se pierde de verdad. La regla que más me servirá es...''}. Firma con nombre y fecha.}

\begin{drawbox}{milcMagenta}{Antes de cerrar la S8}
\checkbox Abrí el historial de versiones de un archivo real. \\[1mm]
\checkbox Identifiqué mínimo 3 versiones con fecha y autor. \\[1mm]
\checkbox Abrí una versión anterior y la comparé con la actual. \\[1mm]
\checkbox Practiqué restaurar una versión y volver a la actual. \\[1mm]
\checkbox Hice boceto del panel de historial en el cuaderno. \\[1mm]
\checkbox Las 4 reglas del uso responsable están firmadas.
\end{drawbox}

\begin{softbox}{milcMagenta}{milcGris}{Plantilla de trabajo}
\textbf{Estructura.} \textit{Paso 1:} escoger archivo. \textit{Paso 2:} abrir historial (clic derecho). \textit{Paso 3:} 3 versiones identificadas. \textit{Paso 4:} comparar versión anterior con actual. \textit{Paso 5:} restaurar y volver. \textit{Paso 6:} boceto + 4 reglas + reflexión + firma.
\end{softbox}

\begin{infoband}{milcMagenta}


\textbf{Lo que va al cuaderno (Actividad 3 · EVALÚA).} Título: «Actividad 3 --- Mi auditoría del historial». \textbf{Formato:} lista de versiones + diferencias detectadas + boceto + reglas firmadas. \textbf{Extensión:} 1 página. \textbf{Sabes que terminaste cuando} confías en que tu trabajo está seguro porque conoces cómo recuperar versiones anteriores.
\end{infoband}

\puente{El producto es la auditoría + dibujo de la pantalla + 4 reglas firmadas.}

\milcestacion{milcMostaza}{Producto de la sesión}
\textbf{Auditoría del historial · 3 versiones identificadas · restauración practicada · 4 reglas firmadas}.
\begin{softbox}{milcMostaza}{milcGris}{Criterios de calidad}
(1) Hay 3 versiones distintas identificadas con fecha y autor. (2) El estudiante comparó al menos 1 versión anterior con la actual. (3) El estudiante practicó restaurar una versión y volver a la original. (4) El boceto del panel de historial está en el cuaderno. (5) Las 4 reglas del uso responsable están firmadas.
\end{softbox}

\puente{Antes de salir, verifica que puedes navegar entre versiones y restaurar una.}

\milcestacion{milcVino}{Evaluación liberadora · cinco dimensiones}

\begin{softbox}{milcVino}{milcGris}{Sentido de la evaluación}
Evaluar no es solo revisar una nota. Es reconocer qué aprendí, qué cambió en mí, qué emociones manejé, cómo me relaciono con la ciudad y la comunidad, y qué le dejaré a quien viene detrás.
\end{softbox}

\textbf{Mi reflexión liberadora en cinco dimensiones.} Completa cada frase iniciada:

\small
\begin{tabularx}{\textwidth}{>{\bfseries\RaggedRight\arraybackslash}p{3.4cm}Y>{\RaggedRight\arraybackslash}p{4.6cm}}
\rowcolor{milcVino}\color{white}Dimensión & \color{white}\bfseries Mi respuesta (frame starter) & \color{white}\bfseries Algo que descubrí\\
1. Personal & Lo que más cambió en mí fue \linead{6.5cm} & \linead{4.2cm}\\
2. Emocional & La emoción más fuerte fue \linead{2.5cm} y la regulé \linead{3cm} & \linead{4.2cm}\\
3. Ciudadana & En democracia esto importa porque \linead{6cm} & \linead{4.2cm}\\
4. Local (Cartago) & En mi barrio sería útil para \linead{6.5cm} & \linead{4.2cm}\\
5. Intergeneracional & A mi abuela/abuelo le diría \linead{6.5cm} & \linead{4.2cm}\\
\end{tabularx}
\normalsize

\begin{infoband}{milcVino}
\textbf{Para la próxima sustentación.} Vuelve a esta tabla en seis meses. Verás que las respuestas cambian — no porque lo aprendido cambie, sino porque tú cambias. La evaluación liberadora es una conversación contigo a través del tiempo.
\end{infoband}


\puente{Cerramos con 3 voces sobre por qué guardar lo que se ha hecho es disciplina del trabajo serio.}

\milcestacion{milcGradeDark}{Triángulo de pensamiento · cierre filosófico}

\begin{softbox}{milcVino}{milcGris}{Dussel · lente del nosotros}
\textit{``Lo que se guarda se puede defender. Lo que se borra sin registro, no.''}

Hay culturas que guardan todo lo que escriben (\emph{archivos de tribunales, actas de juntas, cartas familiares}) y culturas que botan sin registro. \textbf{Las primeras pueden defenderse cuando alguien las acusa injustamente}: \emph{``aquí está el documento que prueba''}. Las segundas dependen de la memoria humana, que falla. OneDrive es un archivo automático que defiende tu trabajo: \emph{``aquí está la versión de hace 2 semanas, prueba que yo sí entregué a tiempo''}. Esa virtud (guardar para defender) es antigua y sigue siendo poderosa.

\textbf{¿Estoy aprovechando el archivo automático que OneDrive me da, o lo ignoro?}
\end{softbox}
\linead{\linewidth}\\[1mm]
\linead{\linewidth}

\begin{softbox}{milcMostaza}{milcGris}{Estoicismo · Marco Aurelio · cuidado interior}
\textit{``Quien registra su trabajo no teme el olvido. Quien no registra, depende de la memoria, que es traidora.''}

Marco Aurelio escribía sus reflexiones diarias en un cuaderno (\emph{Meditaciones}). Por eso las conservamos 1800 años después. Si las hubiera dejado solo en su cabeza, se habrían perdido. \textbf{Tu trabajo en OneDrive es como esas Meditaciones}: queda registrado paso a paso. El historial de versiones es tu cuaderno automático. \emph{No depende de tu memoria. No depende de que recordaras guardar. Está ahí porque el sistema lo hizo por ti}. Esa seguridad permite hacer trabajo más arriesgado y creativo.

\textbf{¿Confío más en mi memoria que en mis archivos, sabiendo que la memoria falla?}
\end{softbox}
\linead{\linewidth}\\[1mm]
\linead{\linewidth}

\begin{softbox}{milcTurquesa}{milcGris}{Floridi · lente de la infoesfera}
\textit{``El historial automático es la promesa de la nube: que tu trabajo no depende de tu propia disciplina perfecta. La nube te perdona errores.''}

Antes, si se te olvidaba guardar con \texttt{Ctrl+G} en Word de escritorio, perdías horas. \textbf{La nube cambió eso: te perdona los olvidos}. Pero la condición es que \emph{conozcas el historial y sepas usarlo}. Quien no sabe, pierde igual que antes. Quien sabe, opera con seguridad. La diferencia entre los dos es 10 minutos de aprender una herramienta. \emph{Hoy fue ese día}.

\textbf{¿Voy a aprovechar el perdón automático de la nube, o seguir trabajando con miedo a perder?}
\end{softbox}
\linead{\linewidth}\\[1mm]
\linead{\linewidth}

\begin{infoband}{milcNegro}
\textbf{Nota para el docente.} Las tres formulaciones del triángulo son la síntesis del autor de la guía, no citas textuales de los pensadores. Si vas a citarlos en clase, remite a la obra original.
\end{infoband}

\milcestacion{milcVino}{Mi autoevaluación · marca una casilla por fila}

{\small
\setlength{\extrarowheight}{2.2pt}
\begin{tabularx}{\textwidth}{>{\RaggedRight\arraybackslash\bfseries}p{3.0cm}YYY}
\rowcolor{milcVino}\color{white}\bfseries Criterio & \color{white}\bfseries Lo logré & \color{white}\bfseries En proceso & \color{white}\bfseries Necesito apoyo\\[.6mm]
Comprensión del tema & \milcopcion{Explico las ideas con mis palabras.} & \milcopcion{Reconozco las ideas, falta precisión.} & \milcopcion{Necesito repasar lo básico.}\\[.8mm]
\rowcolor{milcGris}Pensamiento computacional & \milcopcion{Usé los pilares al armar mi producto.} & \milcopcion{Usé algunos pilares.} & \milcopcion{Necesito releer la estación 2.}\\[.8mm]
Aplicación y producto & \milcopcion{Mi producto cumple los criterios.} & \milcopcion{Está armado, con ajustes pendientes.} & \milcopcion{Aún está en construcción.}\\[.8mm]
\rowcolor{milcGris}Reflexión 5 dimensiones & \milcopcion{Respondí cada una con honestidad.} & \milcopcion{Respondí algunas con detalle.} & \milcopcion{Mis respuestas son muy generales.}\\[.8mm]
Triángulo de pensamiento & \milcopcion{Conecté las 3 voces con mi trabajo.} & \milcopcion{Conecté al menos una.} & \milcopcion{No las relaciono con mi proceso.}\\[.8mm]
\end{tabularx}}

\vspace{3mm}
\textbf{Hoy entendí que:}\\[1mm]
\linead{\linewidth}\\[1mm]
\linead{\linewidth}

\textbf{Mi compromiso para la próxima sesión es:}\\[1mm]
\linead{\linewidth}\\[1mm]
\linead{\linewidth}

\begin{softbox}{milcMostaza}{milcGris}{Cita de despedida · Marco Aurelio}
\textit{``El obstáculo en el camino se vuelve el camino. Lo que estorba al camino se vuelve camino.''}\\[1mm]
Cada error de hoy, bien observado, se convierte en pieza del camino que pisarás mañana.
\end{softbox}

\small
\begin{tabularx}{\textwidth}{>{\bfseries}p{3.6cm}Y}
\rowcolor{milcGradeDark!12}Nombre completo: & \linead{10cm}\\
Fecha: & \linead{4cm} \quad Curso/grupo: \linead{4cm}\\
Mi firma: & \linead{9cm}\\
\end{tabularx}
\normalsize

\begin{infoband}{milcGradeDark}
\textbf{Para la próxima sesión.} Esta semana, abre el historial de versiones de al menos 2 archivos distintos. Familiarízate. Verifica que entiendes cómo navegar entre versiones. La próxima clase aprendes \textbf{Outlook y Teams}: el correo institucional + las reuniones virtuales profesionales. El periodo entra en su recta final.
\end{infoband}

\end{document}
